\section{复环面}
\label{sec:complex-tori}
% ============================================================================

前面我们一直在研究上半平面 $\HH$ 上的函数（模形式）
以及 $\SL_2(\Z)$ 在 $\HH$ 上的作用（基本域、尖点）。
现在要回答一个自然的问题：
$\SL_2(\Z) \backslash \HH$（所有轨道的集合）到底"是"什么空间？

答案是：它参数化了所有的\textbf{复环面}。
复环面是最简单的、非平凡的紧黎曼面，也是椭圆曲线的复分析版本。

\textbf{什么是格？}
在复平面 $\C$ 中，选两个 $\R$-线性无关的复数 $\omega_1, \omega_2$
（即 $\omega_1/\omega_2 \notin \R$），它们生成一个"网格"：
\[
  \Lambda = \Z\omega_1 + \Z\omega_2
  = \{m\omega_1 + n\omega_2 \mid m, n \in \Z\}.
\]

\begin{example}[几个具体的格]
\label{ex:lattices}
\leavevmode
\begin{itemize}
  \item $\Lambda = \Z + \Z i$：正方形格，格点分布在整数坐标上。
    基本平行四边形是单位正方形。
  \item $\Lambda = \Z + \Z e^{2\pi i/3}$：正三角形格（蜂巢状），
    $\omega_2 = -1/2 + i\sqrt{3}/2$。
    基本平行四边形是 $60°$ 的菱形。
  \item $\Lambda = \Z + \Z(2i)$：长方形格，
    基本平行四边形是 $1 \times 2$ 的长方形。
\end{itemize}
\end{example}

\begin{definition}[格与复环面]
\label{def:lattice-torus}
\index{格 (lattice)}\index{复环面 (complex torus)}
$\C$ 中的一个\textbf{格}（lattice）是
$\Lambda = \Z\omega_1 + \Z\omega_2$（$\omega_1/\omega_2 \notin \R$）。

对应的\textbf{复环面}为商群
\[
  \C/\Lambda = \{z + \Lambda \mid z \in \C\}.
\]
两个复数 $z_1, z_2$ 在 $\C/\Lambda$ 中代表同一个点，
当且仅当 $z_1 - z_2 \in \Lambda$（即它们差一个格向量）。
\end{definition}

\textbf{复环面长什么样？}
可以这样想象：取格的一个\textbf{基本平行四边形}
（以 $0, \omega_1, \omega_2, \omega_1+\omega_2$ 为顶点），
然后把对边粘合——上下粘合得到一个圆柱，再把圆柱的两端粘合得到一个甜甜圈（torus）。
这就是为什么叫"环面"。

拓扑上，$\C/\Lambda$ 是一个亏格为 $1$ 的紧曲面。
但它不仅有拓扑结构——从 $\C$ 继承的复结构使它成为一个\textbf{黎曼面}
（Riemann surface，详见\cref{ch:riemann-surfaces}），
而且加法 $z + w$ 使它成为一个\textbf{复 Lie 群}（即加法群）。

\begin{example}[格 $\Z + \Z i$ 对应的环面]
\label{ex:square-torus}
取 $\Lambda = \Z + \Z i$。基本平行四边形是单位正方形
$\{x + iy \mid 0 \leq x < 1,\; 0 \leq y < 1\}$。
\begin{itemize}
  \item 点 $0.3 + 0.7i$ 和 $1.3 + 0.7i$ 在 $\C/\Lambda$ 中是同一个点
    （差 $1 \in \Lambda$）。
  \item 点 $0.5 + 0.5i$ 和 $-0.5 - 0.5i$ 也是同一个点
    （差 $1 + i \in \Lambda$）。
  \item 加法：$(0.6 + 0.8i) + (0.7 + 0.5i) = 1.3 + 1.3i \equiv 0.3 + 0.3i$
    （模掉 $1 + i$）。
\end{itemize}
\end{example}

\begin{proposition}[复环面的同构分类]
\label{prop:torus-isomorphism}
$\C/\Lambda \cong \C/\Lambda'$（作为复 Lie 群）当且仅当
$\Lambda' = \alpha\Lambda$ 对某 $\alpha \in \C^*$。
\end{proposition}

\begin{proof}
设 $\varphi\colon \C/\Lambda \to \C/\Lambda'$ 为全纯群同构。

\textbf{关键思路：从"卷起来的"回到"展开的"。}
复环面 $\C/\Lambda$ 是把复平面 $\C$ 按照格 $\Lambda$ "卷起来"得到的
（就像把一张纸沿两个方向卷成甜甜圈）。
反过来，$\C$ 就是 $\C/\Lambda$ "展开后"的样子——
拓扑学中叫做\textbf{万有覆盖}（universal cover）。
投影映射 $\pi\colon \C \to \C/\Lambda$（$z \mapsto z + \Lambda$）
把 $\C$ 中相差一个格向量 $\omega \in \Lambda$ 的点映到同一个点。

现在，$\varphi$ 是环面之间的映射。我们想知道它"展开后"是什么样的。
覆盖空间理论保证：$\varphi$ 可以"提升"为一个全纯函数
$\tilde\varphi\colon \C \to \C$，使得下图交换：
\[
  \begin{tikzcd}
    \C \arrow[r, "\tilde\varphi"] \arrow[d, "\pi"'] & \C \arrow[d, "\pi'"] \\
    \C/\Lambda \arrow[r, "\varphi"'] & \C/\Lambda'
  \end{tikzcd}
\]
即 $\pi' \circ \tilde\varphi = \varphi \circ \pi$。
直觉上：在展开的平面上做的事情，卷回去之后要和 $\varphi$ 一致。

\textbf{推导 $\tilde\varphi$ 的形状。}
由于 $\pi(z + \omega) = \pi(z)$（差一个格向量不影响），
$\tilde\varphi(z + \omega)$ 和 $\tilde\varphi(z)$ 在 $\C/\Lambda'$
中代表同一个点，即
$\tilde\varphi(z + \omega) - \tilde\varphi(z) \in \Lambda'$。

函数 $z \mapsto \tilde\varphi(z + \omega) - \tilde\varphi(z)$ 是连续的，
但它的值只能落在离散集 $\Lambda'$ 中——
连续函数的像如果是离散的，就必须是常数。
所以对每个固定的 $\omega$，这个差是一个不依赖于 $z$ 的常数。

对 $z$ 求导：$\tilde\varphi'(z + \omega) = \tilde\varphi'(z)$，
即 $\tilde\varphi'$ 以 $\Lambda$ 为周期。
一个全纯的双周期函数（椭圆函数）如果没有极点，
由 Liouville 定理必为常数（\cref{thm:elliptic-basic}(1)）。
所以 $\tilde\varphi'(z) = \alpha$ 是常数，积分得 $\tilde\varphi(z) = \alpha z + \beta$。

最后，$\tilde\varphi$ 把 $\Lambda$ 映到 $\Lambda'$ 内（因为差值在 $\Lambda'$ 中），
又由满射性得 $\alpha\Lambda = \Lambda'$。
\end{proof}

\begin{corollary}
\label{cor:torus-tau}
每个复环面同构于 $\C/\Lambda_\tau$（$\Lambda_\tau = \Z\tau + \Z$，
$\tau \in \HH$）。两个环面 $\C/\Lambda_\tau$ 和 $\C/\Lambda_{\tau'}$
同构当且仅当 $\tau' = \gamma(\tau)$ 对某 $\gamma \in \SL_2(\Z)$。

因此，\textbf{复环面的同构类由 $\SL_2(\Z) \backslash \HH$ 的点参数化}。
这就是上半平面与模群的几何意义——它们在"管理"所有复环面。
\end{corollary}

\begin{proof}
\textbf{第一步：标准化格。}
任意格 $\Lambda = \Z\omega_1 + \Z\omega_2$。
不妨设 $\im(\omega_1/\omega_2) > 0$（否则交换 $\omega_1, \omega_2$）。

乘以 $\alpha = 1/\omega_2$：$\alpha\Lambda = \Z(\omega_1/\omega_2) + \Z$。
令 $\tau = \omega_1/\omega_2$，则 $\alpha\Lambda = \Z\tau + \Z = \Lambda_\tau$，
且 $\tau \in \HH$。
由\cref{prop:torus-isomorphism}，$\C/\Lambda \cong \C/\Lambda_\tau$。

\textbf{关键直觉}：任何格都可以通过缩放变成"第一个生成元在上半平面、
第二个生成元为 $1$"的标准形式。

\textbf{第二步：什么时候两个标准格相同？}
$\C/\Lambda_\tau \cong \C/\Lambda_{\tau'}$ 当且仅当
$\Lambda_{\tau'} = \alpha\Lambda_\tau$ 对某 $\alpha \in \C^*$。

$\Lambda_\tau = \Z\tau + \Z$，所以 $\alpha\Lambda_\tau = \Z(\alpha\tau) + \Z\alpha$。
要求这等于 $\Lambda_{\tau'} = \Z\tau' + \Z$，
即 $\alpha\tau$ 和 $\alpha$ 都是 $\tau'$ 和 $1$ 的整系数线性组合：
\[
  \alpha\tau = a\tau' + b, \qquad \alpha = c\tau' + d,
  \qquad a,b,c,d \in \Z.
\]
消去 $\alpha$：$\tau = \frac{a\tau' + b}{c\tau' + d}$。
由 $\im(\tau) > 0$、$\im(\tau') > 0$ 推出 $ad - bc > 0$。
反过来 $\tau'$ 也要能用 $\tau$ 表示（因为是同构不是单方向），
所以 $ad - bc = 1$，即 $\begin{pmatrix} a & b \\ c & d \end{pmatrix} \in \SL_2(\Z)$。
\end{proof}

\begin{example}[同一个环面的不同参数]
\label{ex:equivalent-tau}
取 $\tau = 2i$ 和 $\tau' = -1/(2i) = i/2$。
它们满足 $\tau' = S(\tau)$（$S \in \SL_2(\Z)$），
所以 $\C/\Lambda_{2i} \cong \C/\Lambda_{i/2}$。
具体地：$\Lambda_{2i} = \Z(2i) + \Z$，
取 $\alpha = 1/(2i) = -i/2$，则
$\alpha \cdot \Lambda_{2i} = \Z(1) + \Z(-i/2) = \Z + \Z(i/2) = \Lambda_{i/2}$。
\end{example}


% ============================================================================
